\subsection{Probability Distributions of Discrete Random Variables}
\hrulefill

\subsubsection*{Probability Distributions}
\paragraph*{Definition}
To describe a discrete random variable X, you need two pieces of information:
\begin{itemize}
    \item What values X can assume.
    \item The probabilities associated with each of these values.
\end{itemize}

The \textbf{probability distribution} or \textbf{probability mass function} (pmf) of a discrete random variable X
is a formula, table, or graph that gives both the possible values of X and the probabilities associated with each value 
of X. The probability distribution or pmf of a discrete random variable is defined for every number x by:

\[
    p(x) = P(X = x) P(\forall s \in S, X(s) = x)
\]

\paragraph*{Example:}
Let X be the random variable that gives the sum of 2d6.
\begin{table}[h!]
    \centering
    \begin{tabular}{c|c|c|c|c|c|c|c|c|c|c|c|c}
        x & 2 & 3 & 4 & 5 & 6 & 7 & 8 & 9 & 10 & 11 & 12 \\
        \hline
        p(x) & $\frac{1}{36}$ & $\frac{2}{36}$ & $\frac{3}{36}$ & $\frac{4}{36}$ & $\frac{5}{36}$ & $\frac{6}{36}$ & $\frac{5}{36}$ & $\frac{4}{36}$ & $\frac{3}{36}$ & $\frac{2}{36}$ & $\frac{1}{36}$
    \end{tabular}
\end{table}

\subsubsection*{Bernoulli Random Variables}
\paragraph*{Definition}
Any random variable whose only possible values are 0 and 1 is called a \textbf{Bernoulli random variable}.
The PMF looks like this:
\begin{table}[h!]
    \centering
    \begin{tabular}{c|c|c}
        x & 0 & 1 \\
        \hline
        p(x) & $1-p$ & $p$
    \end{tabular}
\end{table}

\paragraph*{Example:}
A box contains one red, two orange, and five black balls. You choose two balls at random. For each red ball you win \$5,
for each orange ball you win \$1, and for each black ball you win nothing. Let X be the random variable that gives your 
total winnings. Find the probability distribution of X and the probability that you will win at most \$3.

\begin{table}[h!]
    \centering
    \tabular{|c|c|c|}
        \hline
        x & p(x) & Explanation \\
        \hline
        0 & $\frac{10}{28}$ & BB \\
        \hline
        1 & $\frac{10}{28}$ & BO, OB \\
        \hline
        2 & $\frac{1}{28}$ & OO \\
        \hline
        5 & $\frac{5}{28}$ & RB, BR \\
        \hline
        6 & $\frac{2}{28}$ & RO, OR \\
        \hline
    \endtabular
\end{table}

\[
    P(X \leq 3) = p(0) + p(1) + p(2) = \frac{10}{28} + \frac{10}{28} + \frac{1}{28} = \frac{21}{28} = \frac{3}{4}
\]

\subsubsection*{Cumulative Distribution Function}
\paragraph*{Definition}
The \textbf{cumulative distribution function} (CDF) $F(x)$ of a discrete RV $X$ with PMF $p(x)$ is defined for \textbf{every number x} by:
\[
    F(x) = P(X \leq x) = \sum_{i \leq x} p(i)
\]

\paragraph*{Example:}
Take the previous example with the balls. Some examples of using CDF here are:
\begin{align*}
    F(3) &= P(X \leq 3) = p(0) + p(1) + p(2) = \frac{21}{28} = \frac{3}{4} \\
    F(5) &= P(X \leq 5) = p(0) + p(1) + p(2) + p(5) = \frac{26}{28} = \frac{13}{14} \\
    F(-10) &= P(X \leq -10) = 0 \\
    F(10) &= P(X \leq 10) = 1
\end{align*}

We can find the CDF:
\begin{align*}
    F(x) = 
    \begin{cases}
        0 & x < 0 \\
        \frac{10}{28} & 0 \leq x < 1 \\
        \frac{20}{28} & 1 \leq x < 2 \\
        \frac{21}{28} & 2 \leq x < 5 \\
        \frac{26}{28} & 5 \leq x < 6 \\
        1 & x \geq 6
    \end{cases}
\end{align*}

\paragraph*{Example:}
Find the PMF from the following CDF:
\begin{align*}
    \begin{cases}
    0 & x < -3 \\
    0.2 & -3 \leq x < -2 \\
    0.3 & -2 \leq x < 0 \\
    0.6 & 0 \leq x < 1 \\
    0.9 & 1 \leq x < 4 \\
    1 & x \geq 4
    \end{cases}
\end{align*}

PMF:

\begin{table}[h!]
    \centering
    \begin{tabular}{c|c|c}
        x & p(x) & Explanation \\
        \hline
        -3 & 0.2 & $F(-3) - F(-\infty)$ \\
        \hline
        -2 & 0.1 & $F(-2) - F(-3)$ \\
        \hline
        0 & 0.3 & $F(0) - F(-2)$ \\
        \hline
        1 & 0.3 & $F(1) - F(0)$ \\
        \hline
        4 & 0.1 & $F(4) - F(1)$ \\
        \hline
    \end{tabular}
\end{table}